problem:  
Let \((M^n,g,f)\) be a complete, noncompact **gradient shrinking Ricci soliton** satisfying  
\[
\mathrm{Ric}_g + \nabla^2 f = \tfrac{1}{2}g, \qquad \min f = 0,
\]
and assume \(M\) has **Euclidean volume growth** and **positive curvature operator**.  
It is known that, after suitable rescalings \((M,r_i^{-2}g,p)\), one may obtain an **asymptotic cone** \( (C(Z),g_C) \) with smooth Einstein link \( (Z,h) \), where \( \mathrm{Ric}_h = (n-2)h \).

**Problem:**  
The soliton identity  
\[
R + |\nabla f|^2 - f = C
\]
implies that \(f\) behaves quadratically at infinity, but the exact **higher‑order asymptotics** of \(f\) remain poorly understood.  

1. (**Existence Question**)  
Does every such shrinker admit a **second‑order asymptotic expansion** of the form  
\[
f(x) = \tfrac{1}{4} d(x,p)^2 + \phi(x),
\]
where the remainder \(\phi(x)\) satisfies a well‑defined decay rate relative to the cone structure?  
Equivalently, can one formulate an asymptotic potential correction \(\phi = O(r^{2-\delta})\) for some \(\delta>0\) that remains invariant under blow‑down scaling?  

2. (**Spectral/Geometric Constraint Question**)  
If such a correction \(\phi\) exists, does it satisfy (to leading order on the cone \(C(Z)\)) an elliptic equation of the form  
\[
\mathcal{L}_C \phi = \lambda\,r^{-2} \phi,
\]
where \(\mathcal{L}_C\) is the natural linearization of the shrinking soliton operator on the cone, and \(\lambda\) depends on spectral data of the Laplacian on the Einstein link \(Z\)?  
In particular, can the **spectrum of \(\Delta_Z\)** determine or restrict the possible rates of decay or oscillation of \(\phi\)?  

This formulation asks whether the potential’s deviation from the quadratic leading term carries **analytic fingerprints** of the cone’s link geometry, potentially distinguishing non‑isometric shrinkers with coincident cone limits.

Why is it a "good" problem:  
This version avoids assuming an expansion that has not yet been proved and transforms that assumption into a central question. It targets a fundamental unknown: whether every asymptotically conical shrinking soliton admits a quantitative **second‑order asymptotic expansion of its potential**, and whether that correction encodes **spectral information of the Einstein link**.  
The problem lies at the intersection of PDE analysis (linearization of the soliton equation), spectral theory on cones, and geometric convergence of Ricci solitons—all active, unsolved areas.  
A positive answer would establish a new analytic bridge between soliton geometry and spectral invariants of Einstein manifolds, offering a second‑order theory of asymptotics for shrinkers analogous to the role of ADM mass in asymptotically flat geometry.  
It is conceptually simple, concretely tied to known identities, and yet opens a deep and largely unexplored research frontier in the analytic and geometric structure of Ricci solitons.